\chapter{SecDB and Slang}
\label{ch:secdb_slang}
\chaptermark{SecDB and Slang}

SecDB is Goldman's risk/pricing database; Slang is the proprietary
language that runs inside it.  Together they price a
trillion-dollar derivatives book overnight.  Concretely: SecDB is a
real-time, in-memory graph database holding the firm's entire
position and risk state (every trade, curve, vol, Greek, and
P\&L), updated as bookings and market data arrive, so any desk
can query live risk at any instant.  A CS reader should not picture
SQL: picture a reactive spreadsheet the size of a bank, where every
cell recomputes automatically when any input changes.  SecDB and Slang are
internal to one firm: no vendor manual, no open-source repo.
What we know comes from interviews with former Goldman engineers,
Mike Dubno's InfoQ talk on the Slang VM, and trade-press coverage.
Most claims here are \textbf{training-knowledge-based}, reflecting
the public record as of May 2025; treat the chapter as an
\emph{informed sketch}, correct enough for an interview but
expect details inside the firm to differ.  SecDB is infrastructure,
not a strategy: every Greek in \cref{ch:risk_production}, every
P\&L Explain, every trade ticket lives as a node in a graph
database whose ancestor was written in the 1990s.

\begin{backgroundbox}[Prerequisites]
From earlier chapters:
\begin{itemize}
  \item \Cref{ch:strats_role}: what a Strat does, the divisional
    structure of the Engineering division, and the observation
    that ``Strats code in Slang against SecDB'' is the single
    sentence most often repeated in Goldman job descriptions.
  \item \Cref{ch:risk_production}: portfolio-scale Greeks, P\&L
    Explain, and the idea that pricing a derivatives book is a
    \emph{graph computation}: a curve changes, a million
    sensitivities must follow.  SecDB is the object that makes
    this tractable.
\end{itemize}
This chapter needs no new mathematics.  It is an engineering
chapter about a database schema and a domain-specific language.
\end{backgroundbox}


%% ================================================================
\section{SecDB architecture: a graph that prices a bank}
\label{sec:secdb_arch}
%% ================================================================

Consider Goldman's problem at 5pm New York on a Tuesday.  The firm
holds five million live derivative trades (swaps, swaptions,
exotics, bonds, FX forwards, equity options), each with a
mark-to-market value, Greeks, CVA, FVA.  Those numbers depend on
today's market data: zero curves in every currency, vol surfaces
on every underlying, credit spreads for every counterparty.  By
the time Tokyo opens at 7pm New York, every trade must be
repriced, its Greeks recomputed, and the risk distributed to every
downstream system.

Repricing five million trades from scratch every night is hopeless:
slow, and, more importantly, it gives the Strat no way to
ask ``\emph{which} input moved this number?'', the substance of
P\&L Explain (\cref{ch:risk_production}).  SecDB's answer is to
store the bank as a single \emph{directed acyclic graph} (DAG) of
dependencies: every number is a node, every pricing relationship
an edge.

\subsection{Nodes and edges}

A SecDB \emph{node} is a financial object with a value and a set of
inputs.  Typical nodes:
\begin{itemize}
  \item An individual \textbf{trade}: a five-year USD swap, a
    six-month Brent Call.  Its value depends on curves, vols, spot.
  \item A \textbf{curve}: the USD SOFR zero curve, bootstrapped
    from quoted swap rates.  Its value depends on the market quotes
    that fed the bootstrap.
  \item A \textbf{quote}: the 2y swap rate, a leaf node fed by
    the market-data system.
  \item A \textbf{Greek}: the delta of a specific trade to a
    specific input bucket.  Its value is itself a node, computed by
    bumping one input and re-evaluating the trade.
  \item A \textbf{portfolio aggregate}: the desk's total dollar
    delta to SOFR 2y, summed over every position whose pricer
    depends on that bucket.
\end{itemize}
An \emph{edge} from $A$ to $B$ says ``$B$ depends on $A$''.  A
trade's value node has edges from each curve and vol surface it
consumes; a curve has edges from its bootstrapping quotes; a
portfolio aggregate has edges from every trade it sums.

The graph is a DAG: cycles are refused at commit time, because
they would produce undefined values.

\begin{keyfactbox}[SecDB graph architecture]
\begin{itemize}
  \item Every financial number in the bank (trade value, curve
    point, vol, Greek, aggregate) is a node.
  \item Every pricing dependency is a directed edge.
  \item The whole structure is a DAG; cycles are prohibited.
  \item When a leaf node (e.g.\ a market quote) changes, SecDB
    automatically re-evaluates every descendant node in
    topological order.  No Strat writes ``update everything that
    depends on SOFR 2y''; that is what the graph is \emph{for}.
  \item The graph is persisted: each node is a row in the
    database with a history, in addition to its in-memory value.
\end{itemize}
\end{keyfactbox}

\subsection{Propagation and sensitivities}
\label{sec:secdb_prop}

The central operation is \emph{graph propagation}: a leaf node
changes; SecDB walks the DAG downstream, recomputing each affected
node exactly once in topological order.  Computation is lazy
(recomputed only on query, unless marked eager) and cached.

Graph-native evaluation turns sensitivity calculation from a
bespoke task into a one-liner.  To compute
$\partial V / \partial r$ for a trade whose value $V$ depends on a
rate input $r$:
\begin{enumerate}
  \item Bump the leaf node $r \to r + \epsilon$.
  \item Ask SecDB for $V$ (``evaluate node $V$'').
  \item SecDB propagates $r+\epsilon$ downstream, recomputes every
    dependent node, and returns the new $V$.
  \item Sensitivity $= (V_{\text{new}} - V_{\text{base}}) / \epsilon$.
\end{enumerate}
The Strat writes no plumbing: she does not enumerate which curves
rebuild, which surfaces refit, which trades reprice.  The graph
does it.  This is SecDB's most important property:
\emph{sensitivities are first-class}.  A book-level
dollar-delta-to-SOFR-2y has the same structure as a single-trade
dollar delta; only the downstream nodes differ.

The naive approach (``bump one factor, reprice from scratch,
repeat $N$ times'') scales as $O(N \cdot M)$ where $M$ is the
work to price the portfolio.  SecDB's cached graph reprices only
the nodes actually affected, a 10--100x speed-up on a large book.

\begin{intuitionbox}
A spreadsheet with formulas is a graph of cells; change an input
cell, every formula downstream updates.  SecDB is that spreadsheet
scaled to fifty million rows, with the cells persisted to a
database, the formulas written in a specialised language, and the
recalculation engine optimised to remember what it has already
computed.  Everything else in this chapter is engineering
consequences of that one idea.
\end{intuitionbox}

\subsection{Bi-temporal data}
\label{sec:bitemporal}

A plain graph database is insufficient: banks must answer two
kinds of historical question, with different answers:
\begin{enumerate}
  \item ``What was the trade's value \emph{as of} 31 December
    2023?'', the question an auditor asks.  This is
    \textbf{business date} (also called \emph{valuation date} or
    \emph{as-of date}).
  \item ``What was the value of that trade \emph{as we believed
    it} on 2 January 2024 at 9am?'', the question a regulator
    asks after a fat-finger, or a Strat asks after a bug is fixed.
    This is \textbf{system date} (also called \emph{knowledge date}
    or \emph{transaction time}).
\end{enumerate}
The two answers differ whenever a curve is re-bootstrapped, a
trade amendment is booked with a past effective date, or a
pricing-model bug is corrected.  Business date is when the value
\emph{applies}; system date is when it was \emph{entered}.

SecDB is \emph{bi-temporal}: every node carries both timestamps,
and every query specifies both.  The default is ``business =
today, system = now'', but an auditor can replay the graph as it
looked on any historical system date, and a Strat debugging
yesterday's P\&L Explain can ask for the curves as they were fed
into the batch at 6pm last night, not as they look now after an
overnight correction.

\begin{keyfactbox}[Bi-temporal storage]
Every node has two timestamps:
\begin{itemize}
  \item \textbf{Business date}: the date the value refers to
    (``close-of-business 15 Mar 2024'').
  \item \textbf{System date}: the date the value was recorded in
    SecDB (``15 Mar 2024 18:32:04'').
\end{itemize}
A correction booked on 18 March, effective 15 March, creates a new
row with business date 15 March, system date 18 March.  The
original row is not deleted; both are queryable.  Bi-temporal
storage is what makes SecDB \emph{auditable}: every number ever
reported to a regulator can be reconstructed exactly.
\end{keyfactbox}

SecDB grows monotonically (nothing is overwritten), but the
reproducibility is non-negotiable for a regulated bank.

\subsection{A worked bi-temporal example}
\label{sec:bitemporal_example}

Consider a five-year USD interest-rate swap booked on 15 March.
On the evening of the 15th, SOFR 2y fixes at 4.82\% and the
overnight batch marks the trade at \$123{,}456.  The trade's value
node is persisted with:
\[
  (\text{business date} = \text{15-Mar},\;
   \text{system date} = \text{15-Mar 18:32}),\qquad
  V = \$123{,}456.
\]
On the morning of the 18th a correction comes in: the 15th's SOFR
2y fix was republished as 4.79\% by the reference-rate provider.
The batch reruns as of business date 15-Mar, pulling the corrected
rate, and produces a new value \$124{,}010.  SecDB writes:
\[
  (\text{business date} = \text{15-Mar},\;
   \text{system date} = \text{18-Mar 08:14}),\qquad
  V' = \$124{,}010.
\]
The original row is \emph{not} deleted.  Three queries now behave
as follows:
\begin{itemize}
  \item ``Value on 15-Mar as we knew it on the 15th'': returns
    \$123{,}456.  This is the number the trader saw in her Monday
    morning report.
  \item ``Value on 15-Mar as we know it now'': returns \$124{,}010.
    This is the corrected number, the one that appears in the
    audit trail.
  \item ``P\&L Explain for 18-Mar'': reports a \$554 ``market-data
    correction'' line, explaining why the mark moved without any
    real market action.
\end{itemize}
If an auditor next year asks ``what did Goldman believe this
trade was worth on 16-Mar?'', the query (business = 15-Mar,
system = 16-Mar) returns \$123{,}456, the value used in that
day's regulatory capital report.  Both answers are correct, both
retrievable, neither overwrites the other.  That is what
bi-temporal storage \emph{does}.


%% ================================================================
\section{Slang as a language}
\label{sec:slang}
%% ================================================================

Slang is the DSL that runs inside SecDB.  It was designed in the
1990s by Mike Dubno and a small team around the question ``what
does the Strat need to write a pricer, and what can the VM do to
make that writing productive?''  The result is unlike any
mainstream production language.  The mental model is that Slang is
to SecDB what a cell formula is to a spreadsheet: the language in
which pricing rules (option pricers, curve builds, Greek
formulas) are expressed, executing in the same process as the risk
data rather than in a separate client talking to a database.

\subsection{Design choices}

\paragraph{Dynamic typing, C-like syntax.}  Slang looks
superficially like C (curly braces, semicolons, function
declarations) but is dynamically typed: variables hold whatever
they were last assigned; the VM checks types at use time.
Exploratory pricing code writes quickly; refactoring a
fifty-million-line codebase is harder because type errors appear
at runtime.

\paragraph{No keywords: \texttt{if} and \texttt{while} are
functions.}  Control-flow constructs are not built-in keywords
but user-callable functions, and therefore \emph{overridable}.  A
Strat wanting a specialised \texttt{if} that logs both branches
while evaluating only the taken one writes
\texttt{my\_if(cond, t, f)\{...\}}.  Scheme does this; Python does
not.  The design reflects Goldman culture: give Strats the
toolkit and trust them.

\paragraph{Spaced variable names.}  Slang allows spaces in
identifiers: \texttt{average annual forward rate} rather than
\texttt{averageAnnualForwardRate}.  Code reads like financial
quantities.  The parser resolves ambiguity via context; the cost
is that some mainstream constructs (juxtaposition for function
application) are unavailable.

\paragraph{Functional style, limited side effects.}  Slang
encourages pure functions.  The graph engine caches results
aggressively; a function with hidden state would return an
incorrect cached output next time the same inputs appear.  Pure
functions compose cleanly with the cache.  Side-effects exist
(logging, persistence) but are boxed off.

\paragraph{Graph-aware VM.}  The Slang VM knows about SecDB:
function-call results are keyed by the set of graph nodes the
call read.  Same call, unchanged nodes $\to$ cache hit.  A changed
node invalidates every cache entry whose key-set included it.
This is how the bank reprices overnight: most results are already
cached from last night; only the subset whose inputs moved
recomputes.

\begin{keyfactbox}[Slang language features]
\begin{itemize}
  \item Dynamic typing, C-like syntax.
  \item Control-flow constructs (\texttt{if}, \texttt{while},
    \texttt{for}) are functions, not keywords, and are
    overridable.
  \item Identifiers can contain spaces:
    \texttt{forward swap rate} is a legal variable name.
  \item Functional style: pure functions the VM can cache
    correctly.
  \item The Slang VM is graph-aware: cached results are keyed by
    the set of SecDB nodes read, and invalidated when any of those
    nodes changes.
\end{itemize}
\end{keyfactbox}

\subsection{A stylised example}

The following is a stylised Slang-like pricer: pseudocode
capturing the flavour, not real Slang syntax.

\begin{lstlisting}[language=C, basicstyle=\ttfamily\small, frame=single, caption={Stylised Slang-like pricer for a European call (pseudocode)}]
// Pure function: computes BS call given spot, strike,
// rate, vol, and maturity.  Graph-aware: the VM caches
// the result keyed by the five inputs.

black scholes call(spot, strike, rate, vol, maturity) {

    // d1 and d2 under the standard BS parameterisation.
    d1 = (log(spot / strike)
          + (rate + 0.5 * vol * vol) * maturity)
         / (vol * sqrt(maturity));
    d2 = d1 - vol * sqrt(maturity);

    // Normal CDF as an ordinary function call.
    return spot * norm cdf(d1)
         - strike * exp(-rate * maturity) * norm cdf(d2);
}

// Trade pricer reads SecDB nodes for each input and
// delegates to the BS pricer.  Because the nodes are
// inputs to a pure function, the VM knows which nodes
// this result depends on and will invalidate correctly.

price european call trade(trade) {
    spot     = market spot[trade.underlying];
    strike   = trade.strike;
    rate     = discount curve[trade.ccy].rate(trade.maturity);
    vol      = vol surface[trade.underlying]
                 .implied vol(trade.strike, trade.maturity);
    maturity = year fraction(today, trade.maturity);
    return black scholes call(spot, strike, rate,
                              vol, maturity);
}
\end{lstlisting}

A second snippet shows the purity \emph{contract} for safe VM
caching, contrasted with a subtly-wrong version:

\begin{lstlisting}[language=C, basicstyle=\ttfamily\small, frame=single, caption={Pure vs impure: why the VM cache depends on purity}]
// CORRECT: pure, graph-aware.  All reads go through
// declared SecDB inputs; no hidden state.

pv of swap(trade, valuation date, discount curve) {
    legs = generate cashflows(trade, valuation date);
    return sum over(legs,
                    leg -> discount curve.df(leg.pay date)
                         * leg.amount);
}

// WRONG: reads today from a global mutable clock.  The VM
// caches the result keyed only by (trade, valuation date,
// discount curve), but today is a hidden input.  Tomorrow
// the cache returns yesterday's answer.

pv of swap bad(trade, discount curve) {
    legs = generate cashflows(trade, global today());
    return sum over(legs,
                    leg -> discount curve.df(leg.pay date)
                         * leg.amount);
}
\end{lstlisting}

The fix: promote the hidden input to an explicit argument, and
let the VM track it as a graph node.  The Slang style guide
reduces to ``make your functions pure or you will poison the
cache.''

Three features of real Slang show up in the stylised form: spaced
identifiers (\texttt{black scholes call}); function composition
whose inputs are all SecDB nodes, so the VM tracks exactly which
cache entries to invalidate; and no explicit types, all
inferred from usage.

\subsection{The VM's job: aggressive caching}

Most of Slang's cleverness is in the VM, not the language surface.
Without caching, overnight repricing would re-execute every
function in every trade.  With caching, the VM observes:
\begin{itemize}
  \item \texttt{black scholes call(100, 100, 0.05, 0.2, 1.0)} has
    already been computed once today, returning 10.45.  If
    queried again with identical inputs, the VM returns the cache.
  \item \texttt{price european call trade(trade\_42)} previously
    depended on \texttt{market spot["AAPL"]},
    \texttt{discount curve["USD"]}, and
    \texttt{vol surface["AAPL"]}.  The SOFR 2y bucket moved;
    \texttt{discount curve["USD"]} invalidates;
    \texttt{price european call trade(trade\_42)} invalidates;
    all dependents of that trade invalidate.
\end{itemize}
Memoisation at whole-language scale: every pure call is memoised,
and invalidation is driven by graph edges, not annotations.  The
overnight batch recomputes only the numbers the market moved.

\subsection{A day in the life of a graph change}
\label{sec:day_in_life}

A concrete scenario.  At 5:02pm New York on Tuesday, the
SOFR-curve build has just run against the 5pm swap-rate snap, and
the 2y-bucket discount factor moved from 0.9100 to 0.9095.  This
single change triggers the following cascade, orchestrated by
SecDB without Strat intervention:
\begin{enumerate}
  \item The node \texttt{usd zero curve.discount(2y)} is marked
    dirty.
  \item Every pricing function whose cached result read that node
    is invalidated.  Call the invalidation set $\mathcal{I}_1$.
    Typically this includes every USD swap and swaption with a
    2y-region sensitivity, every USD bond, every USD-funded exotic
    whose funding leg crosses the 2y bucket.
  \item Every aggregate whose cached result read any node in
    $\mathcal{I}_1$ is marked dirty.  The desk-level dollar-delta
    to SOFR 2y, the firm-level VaR, the USD-book CVA all enter a
    new invalidation set $\mathcal{I}_2$.
  \item A batch scheduler walks the invalidation sets in
    topological order, dispatching recomputations across a
    compute cluster.  Trades in $\mathcal{I}_1$ reprice; then
    aggregates in $\mathcal{I}_2$ are re-aggregated from the new
    trade values; then any reports, dashboards, or regulatory feeds
    that consume those aggregates rebuild.
  \item By 5:07pm the downstream effects have propagated through
    perhaps a hundred thousand nodes; every number on every screen
    in the firm that depends on SOFR 2y now reflects the new
    curve.  The Strat is not involved.
\end{enumerate}
Nobody writes ``if SOFR 2y moves, recompute the following ninety
thousand trades.''  The graph \emph{is} that specification; the
propagation engine reads and acts.


%% ================================================================
\section{Java and Python: the long migration}
\label{sec:migration}
%% ================================================================

By the mid-2010s Slang's drawbacks were worsening with scale.
Dynamic typing made refactoring painful across fifty million
lines.  The hiring pool was exactly one firm.  New graduates came
fluent in Python/Java, and the Slang transition cost months.
Tooling lagged commercial ecosystems by a decade.

Goldman's response was gradual.  Post-2015, \emph{new} development
moves to Java (performance-sensitive production) and Python
(research, scripting), while the SecDB core (graph engine,
persistence, the decades of production pricers) stays in Slang.
The strategy is coexistence: Python and Java read/write SecDB
nodes via bindings; heavy graph computation runs on the Slang VM.

Three practical shapes this takes:
\begin{itemize}
  \item \textbf{Slang $\leftrightarrow$ Python bindings.}  Research
    Strats pull a curve from SecDB, manipulate it in NumPy/pandas,
    and write the result back.  The Python code is not graph-aware
    (dependencies are not tracked for auto-invalidation), but
    for ad-hoc analysis that is acceptable.
  \item \textbf{Java services calling the graph.}  Trading systems,
    risk dashboards, and regulatory reporting increasingly run as
    Java services that query SecDB.  They are plumbing, not
    pricers.  Java is chosen for ecosystem, performance, and
    ubiquity.
  \item \textbf{Slang stays where graph semantics matter.}  A new
    exotic pricer, curve-build function, or xVA calculation (any
    code that must participate in overnight propagation with
    correct cache invalidation) stays in Slang.  The VM's graph
    awareness is not cheaply replicated elsewhere.
\end{itemize}
\subsection{Interop mechanics}

Three interop patterns recur:
\begin{itemize}
  \item \textbf{Pull a node into Python.}  A binding reads a SecDB
    node as-of a (business, system) date pair and materialises it
    as a Python object: NumPy array for a curve, DataFrame for
    a blotter, scalar for a Greek.  The value is a snapshot; if
    the underlying node moves, the Python copy goes stale
    silently.  Fine for notebooks, catastrophic for production
    risk.
  \item \textbf{Write a node back from Python.}  A derived curve
    or scenario written back enters the graph as a new node,
    versioned bi-temporally and attributed to the author.  Writers
    face higher review standards than readers.
  \item \textbf{Call Slang from Java.}  A Java service calls the
    Slang VM via RPC; the Slang side runs the pricer with full
    graph semantics including cache, and returns a value plus
    node-identity for later invalidation queries.  Most Java
    services are thin clients, not pricing implementations.
\end{itemize}

\begin{pitfallbox}[Common mistake: assuming Python owns the state]
A research Strat pulls a SOFR curve into Python, derives a stress
scenario, runs backtests for a week, reports to the desk head.
Between day 1 and day 7, SecDB corrected its 15-March SOFR 2y fix
(\cref{sec:bitemporal_example}).  The Python copy still holds the
\emph{old} value (the pull was a snapshot), so the backtest
silently refers to retracted market data.

The fix is discipline: tag every snapshot with its (business,
system) pair, treat it as read-only, refresh explicitly when
analysis crosses days.  The graph engine tracks staleness for
Slang; Python does not get that for free.
\end{pitfallbox}

Industry reporting calls the migration ``decades-long''.  A Strat
joining Goldman in the mid-2020s meets an actively-maintained
Slang codebase alongside growing Java and Python ecosystems, not
a dying language.  Slang, Python, and Java are all table stakes;
the real skill is knowing \emph{which} codebase to reach for.


%% ================================================================
\section{Shipping code into SecDB}
\label{sec:shipping}
%% ================================================================

A Strat writing a new pricer does not ``merge to master and
deploy.''  SecDB is a shared production database holding a
trillion dollars of risk; any graph change potentially affects
every dependent trade.

\subsection{Environments}

Slang code moves through a sequence of environments:
\begin{enumerate}
  \item \textbf{Local / sandbox.}  The Strat develops against a
    SecDB snapshot; no production trades affected.
  \item \textbf{Staging.}  Promoted to a shared environment where
    QA, tests, and other desk Strats exercise the code against
    realistic data and compare outputs to expected values.
  \item \textbf{Production.}  After review and sign-off, the
    function joins the overnight batch; every dependent trade
    reprices on the new version.
\end{enumerate}
Staging is where bi-temporal storage pays off: replay yesterday's
batch with the new pricer substituted in, look at the P\&L diff,
and demonstrate the economic impact before touching production.

\subsection{Versioning and auditability}

Every Slang function change is versioned bi-temporally in SecDB
itself: Slang code lives in the same graph it computes on.
When a function updates:
\begin{itemize}
  \item The previous version is preserved and tagged with its
    business-date range.  Historical reproductions use the old
    version; new calculations use the new one.
  \item Reviewer identity, commit message, and diff are persisted;
    auditors walk back every pricer change in the bank's history.
  \item Fixing a bug does not erase its footprint: the old
    pricer's output remains retrievable as-of any past system
    date, and the correction shows up explicitly in P\&L Explain.
\end{itemize}
Code and data are versioned the same way, so ``show me what
happened on this date'' extends to ``show me what \emph{code} was
running''.

\subsection{Testing and graph-replay}
\label{sec:testing}

Testing a Slang pricer is unusual: the test harness is the same
graph engine as production.  Typical patterns:
\begin{itemize}
  \item \textbf{Snapshot test.}  Run the production and candidate
    pricers on a representative trade set against yesterday's
    business-date data; flag any price diff above threshold.
    Bi-temporal storage makes this reproducible: yesterday's
    data is frozen and queryable.
  \item \textbf{Bump-parity test.}  Finite-difference Greeks
    should agree to leading order with any analytical Greeks the
    pricer exposes.  Disagreement signals an analytical bug or a
    pricer discontinuity the bump stepped across.
  \item \textbf{Invalidation test.}  The CI harness swaps in a
    mutable-global version of a function (the WRONG pattern
    above) and verifies the framework catches cache-pollution.
    This protects the purity contract the graph engine assumes.
  \item \textbf{Shadow run.}  For major changes, run new and old
    pricers in parallel for days or weeks, reporting both into a
    P\&L-diff feed.  When cumulative drift is understood, promote
    and retire.  Shadow runs are the industry's answer to
    ``change an airplane's engine in flight''.
\end{itemize}
Each pattern leans on graph-replay (rerun a calculation with
business date fixed, system date stepped), which falls out of
the same bi-temporal persistence that supports the audit.

\begin{prosperitybox}
Prosperity's \texttt{traderData} JSON is a miniature SecDB.  Each
tick, the simulator hands your algorithm the previous
\texttt{traderData} plus fresh market state and expects an updated
string back.  The string is persistent, versioned (each tick is a
system-date snapshot), and readable after the round.  It is the
only state a Prosperity algorithm controls.

Three SecDB analogies:
\begin{itemize}
  \item Treat \texttt{traderData} as a versioned state graph
    (rolling means, positions, signal caches, regime flags).
    Schema-version the JSON (\texttt{"schema": "v3"}) so you can
    evolve without breaking mid-round.
  \item Cache expensive computations (OU fits, vol estimates)
    keyed by inputs; update only on material moves.  SecDB-style
    memoisation, a thousand times simpler.
  \item Debug post-hoc by replaying ticks with the stored
    \texttt{traderData}, the bi-temporal-replay analogue.
\end{itemize}
The scale is laughable (kilobytes vs.\ petabytes) but the
\emph{shape} (persistent, versioned, auditable state) is
the same.
\end{prosperitybox}

\begin{gsbox}
Four talking points for an interview that turns to SecDB/Slang.
The interviewer is typically an ex-Strat or a Java/Python-forward
Engineer; either way, the bar is ``understands SecDB is a graph
and Slang exploits it'', not ``recites the language spec''.
\begin{itemize}
  \item \emph{Memoisation and invalidation.}  ``The VM caches pure
    functions by the set of SecDB nodes they read.  When a node
    changes bi-temporally (a new system-date entry for the same
    business date), every cache entry whose read-set includes
    that node invalidates, and dependents recompute on next
    access.  The correctness hinges on functions being pure; a
    function with hidden mutable state would return stale cached
    results.''
  \item \emph{Circular dependencies.}  ``The graph is a DAG;
    cycles are refused at commit time.  A subtle cycle case:
    calibrating a vol surface to option prices that themselves
    depend on the surface.  You break the cycle by making the
    calibration a batch step that runs against yesterday's
    surface and outputs today's, rather than a same-date
    recursion.''
  \item \emph{Bi-temporal debugging.}  ``To debug a P\&L Explain
    discrepancy, replay the trade's valuation with business date
    fixed at yesterday and system date walked forward through the
    overnight corrections; the point at which the value jumps
    isolates the corrected input.  This is the Slang debugger's
    most-used feature.''
  \item \emph{Why not rewrite in Python?}  ``Because the graph
    engine, the cache, the bi-temporal persistence, and the
    auditability are not a language feature; they are a
    fifty-million-line ecosystem, and the VM's integration with
    the database is load-bearing.  Python and Java are layered on
    top; Slang is not going away.''
\end{itemize}
\end{gsbox}

%% ================================================================
\section{What the competition looks like}
\label{sec:competition}
%% ================================================================

Every major sell-side bank faces the same problem (price a
huge derivatives book, propagate sensitivities, reproduce
historical valuations), and the responses converge on the same
architectural idea with different implementations.  A Strat
interviewing at Goldman benefits from knowing how the firm's
stack compares.

\subsection{Peer systems}

\textbf{JPMorgan's Athena} is the most-cited peer.  Built by a
team including ex-Goldman engineers, Athena adopts SecDB-like
ideas (graph dependencies, bi-temporal storage) but uses
\emph{Python} as its in-system language.  The design bet: Python's
ecosystem productivity beats a DSL's graph-awareness advantage,
and framework code can recover most of the missing graph
semantics.  Athena is a plausible counterexample to the claim
that SecDB-style systems \emph{require} a bespoke language.

\textbf{Bank of America Merrill's Quartz} and \textbf{Barclays's
BARX / Jupiter} fall into the same family: Python-centric,
graph-structured, bi-temporal.  The common pattern across
Athena/Quartz/Jupiter is that the language choice is the single
biggest divergence from SecDB; the deeper architectural ideas are
the same.

\textbf{Deutsche Bank's dbAnalytics and Morgan Stanley's Merlin}
occupy a middle ground: partially proprietary, partially
off-the-shelf, with various tiers of graph-awareness.  The details
matter less than the generalisation: every large sell-side
derivatives desk runs on a system that you can describe, to first
approximation, as ``graph database plus pricer library plus
bi-temporal persistence.''

\subsection{Why the industry converged here}

Why graph-plus-bi-temporal rather than relational-plus-batch?
Three forcing functions:
\begin{itemize}
  \item \textbf{Sensitivity requests are ad-hoc.}  ``What is my
    P\&L if EUR 10y spreads 1bp?''  A relational system needs
    a purpose-built job per question; a graph system bumps one
    node and queries another.
  \item \textbf{Regulators audit history.}  Bi-temporal storage is
    a prerequisite for Basel and SEC reconstruction, not a
    nice-to-have.  Firms that bolted audit onto single-timestamp
    databases regretted it.
  \item \textbf{Pricer reuse is the biggest development lever.}
    Pricing the same cashflow-discount logic in ten systems costs
    ten times and produces ten answers.  A single graph-aware
    pricer store is a one-time investment that pays back every
    year in consistency.
\end{itemize}
Every sell-side derivatives business ends up at a SecDB-shaped
solution, whether or not the firm calls it that.

\subsection{Why SecDB is distinctive}

Three things set SecDB/Slang apart.  \emph{Language}: no other
major bank kept a proprietary DSL at Goldman's scale; productivity
arguments pushed peers to Python or Java.  \emph{Coverage}: SecDB
prices every derivative across every asset class, an unusual
centralisation; some banks run separate systems per asset
class.  \emph{Age}: the graph design was in place by the
mid-1990s, before commercial vendors had anything comparable, and
Goldman has run a proprietary graph database for three decades
while peers moved to vendor stacks, a cultural choice as much
as a technical one.

\begin{intuitionbox}
A spreadsheet shows the core idea Goldman generalised to bank
scale: cells depend on other cells via formulas, and updates
propagate automatically.  Everything SecDB adds (persistence,
bi-temporality, graph-aware caching, purity-focused language)
is engineering to make the spreadsheet idea survive five million
trades and a regulator.  The interview takeaway: Goldman invested
early in a graph-database model the industry later converged on,
and stayed with the language-level integration others backed
away from.
\end{intuitionbox}


\begin{historybox}[title={SecDB, Slang, and the J. Aron heritage}]
SecDB predates Slang; its origins trace to J. Aron \& Company, a
commodity-trading firm Goldman acquired in 1981.  J. Aron's
traders needed a system to track commodity positions, prices, and
P\&L across a heterogeneous book, and early SecDB grew out of
that tool.  The graph-of-financial-objects emphasis reflects
trading-desk pragmatism: built by people who needed risk answers
in seconds.

Slang was designed in the 1990s by Mike Dubno and a small team
to replace an earlier SecDB-pricer language.  Dubno's InfoQ talk
remains the best public description of the VM.  Three decades on,
both remain in active production: Strats joining Engineering
today commit code to systems whose DNA traces to a 1980s
commodity-trading desk.
\end{historybox}


%% ================================================================
\section{Key facts to memorise}
\label{sec:ch36_key}
%% ================================================================

\begin{keyfactbox}[Chapter summary]
\begin{itemize}
  \item \textbf{SecDB} is Goldman's proprietary
    risk/pricing database; \textbf{Slang} is the proprietary
    language that runs inside it.  Together they price the
    firm's trillion-dollar derivatives book overnight.
  \item SecDB stores the bank as a DAG.  Every financial
    number (trade value, curve, vol, Greek, aggregate) is a
    node; every pricing dependency is an edge.
  \item A leaf node changes; SecDB propagates downstream in
    topological order, recomputing exactly the affected
    descendants.  This graph-native evaluation is why
    sensitivities are a one-liner and overnight batches are
    feasible.
  \item Storage is \emph{bi-temporal}: every node has a business
    date (the date the value applies to) and a system date (the
    date the value was entered).  Corrections are new rows, not
    overwrites.  Bi-temporal storage is the basis of
    reproducibility and audit.
  \item Slang is dynamically typed, C-like in syntax, but has
    no keywords: \texttt{if} and \texttt{while} are functions
    and overridable.  Identifiers can contain spaces.  Functions
    are expected to be pure; the Slang VM caches their results
    keyed by the set of SecDB nodes read, and invalidates via the
    graph.
  \item Post-2015 Goldman direction: new development in Java
    (services) and Python (research), with Slang retained for
    core pricing logic.  The migration is explicitly decades-long;
    SecDB itself stays Slang.
  \item Shipping Slang goes local $\to$ staging $\to$ production,
    with every change versioned bi-temporally in SecDB itself.
    Code and data live in the same graph; both are audit-replayable
    to any historical system date.
  \item Prosperity's \texttt{traderData} JSON is a miniature
    SecDB: persistent, versioned state-passing at a
    kilobyte-scale analogue of the bank's petabyte graph.
  \item Everything in this chapter is
    training-knowledge-based.  Public documentation of SecDB
    and Slang is fragmentary; expect details inside the firm to
    differ.
\end{itemize}
\end{keyfactbox}
